\documentclass[11pt]{article}
\usepackage[margin=1in]{geometry}
\usepackage{amsmath}
\usepackage{hyperref}

\newcommand{\code}{https://github.com/krash3125/cs422-assignment-4/blob/main/topology.py}

\title{CS 422 Assignment 4 \\ Part 2: Arbitrary Traffic Matrix}
\author{Kush, Joseph, Issac}
\date{}

\begin{document}
\maketitle

\section*{Main idea}

In Part 1 every pair sent the same amount, so the best I could do was make every node close to
every other node. Here the traffic matrix $T$ is given, and some pairs send a lot more than others.
So the topology should depend on $T$: pairs with a lot of traffic should get direct links (even more
than one link), and pairs with little traffic can share capacity over 2 or more hops.

Choosing the links and routing the flow depend on each other, so I put both into one integer
program. For each given $T$, I solve it, and the solution gives the links to build. Since the solver
looks at every possible topology and routing and returns the best one, the topology it gives is the
one that maximizes concurrent flow for that $T$. Here $n = 8$ and $d = 4$.

\section*{Formulation}

\paragraph{Decision variables.}
\begin{itemize}
  \item $x_{uv} \in \{0, 1, \dots, 4\}$ (integer), for $u \neq v$: number of links from $u$ to $v$.
        More than one link between the same pair is allowed.
  \item $f^s_{uv} \ge 0$: amount of flow from source $s$ that goes over link $u \rightarrow v$.
  \item $\theta \ge 0$: the fraction of every demand that gets delivered.
\end{itemize}

\paragraph{Objective.}
\[
\max \ \theta
\]

\paragraph{Constraints.}
\begin{enumerate}
  \item \textbf{Degree:} every node has 4 links out and 4 links in.
  \[
  \sum_{v \neq u} x_{uv} = 4, \qquad \sum_{v \neq u} x_{vu} = 4 \qquad \forall u
  \]
  \item \textbf{Flow conservation:} source $s$ sends $\theta \cdot T_{st}$ to every destination $t$.
  \[
  \sum_{v} f^s_{sv} - \sum_{v} f^s_{vs} = \theta \sum_{t} T_{st} \qquad \forall s
  \]
  \[
  \sum_{v} f^s_{vw} - \sum_{v} f^s_{wv} = \theta \, T_{sw} \qquad \forall s, \ \forall w \neq s
  \]
  \item \textbf{Capacity:} each link has capacity 1, so the total flow on $u \rightarrow v$ is at
        most the number of links there.
  \[
  \sum_{s} f^s_{uv} \le x_{uv} \qquad \forall u \neq v
  \]
\end{enumerate}

I group flow by source instead of by (source, destination) pair. The flow from $s$ only needs to
drop off $\theta T_{st}$ at each $t$, so this is enough and uses fewer variables.

\section*{Why this works}

\begin{itemize}
  \item The degree constraints make sure $x$ is a valid topology (4 links in and out per node).
  \item The flow and capacity constraints make sure the traffic $\theta T$ can actually be routed on
        that topology.
  \item The solver maximizes $\theta$ over all valid topologies and routings together, so the answer
        is the best topology for this $T$. It adapts to $T$ because $T$ shows up in the flow
        conservation constraints, so a different $T$ gives different links.
  \item The intuition shows up in the results: big entries of $T$ get direct links, and
        if $T$ is something like $4 \times$ a permutation, the solver just builds those links and gets
        $\theta = 1$.
\end{itemize}

\section*{Implementation (Gurobi)}

Code: \url{\code}
\begin{itemize}
  \item \href{\code\#L11-L49}{\texttt{best\_topology(T)}}: builds and solves the integer program
        and returns $\theta$ and the links.
  \item \href{\code\#L18-L23}{Variables} $x$, $f$, $\theta$;
        \href{\code\#L25}{objective};
        \href{\code\#L27-L30}{degree constraints};
        \href{\code\#L32-L40}{flow conservation};
        \href{\code\#L42-L44}{capacity}.
  \item \href{\code\#L52-L63}{\texttt{random\_hose\_tm()}}: makes a random hose traffic matrix (a random
        mix of permutations, scaled by 4) for testing.
\end{itemize}

\section*{Results}

\begin{center}
\begin{tabular}{l|c}
traffic matrix & best $\theta$ \\ \hline
uniform ($\frac{4}{7}$ to every pair) & 0.700 \\
random hose \#0 & 0.875 \\
random hose \#1 & 0.861 \\
random hose \#2 & 0.827 \\
\end{tabular}
\end{center}

For the uniform matrix the solver gets $0.700$. This matches Part 1: with $d = 4$, each node can
reach 4 nodes in 1 hop and the other 3 in 2 hops, so the average path length is $10/7$ and
$\theta \le 7/10$. The random matrices get higher $\theta$ because their traffic is concentrated on
fewer pairs, so more of it can go over direct links.

\end{document}
